% Generated draft: requires mathematical and source review.
\documentclass[11pt]{article}
\usepackage[margin=25mm]{geometry}
\usepackage{fontspec}
\setmainfont{Latin Modern Roman}
\usepackage{amsmath,amssymb,mathtools,unicode-math}
\setmathfont{Latin Modern Math}
\usepackage{enumitem,needspace,xurl,hyperref}
\setlength{\parindent}{0pt}
\setlength{\parskip}{5pt}
\setlength{\emergencystretch}{3em}
\newcommand{\N}{\mathbb{N}}
\newcommand{\Z}{\mathbb{Z}}
\newcommand{\Q}{\mathbb{Q}}
\newcommand{\R}{\mathbb{R}}
\newcommand{\C}{\mathbb{C}}
\newcommand{\F}{\mathbb{F}}
\newcommand{\A}{\mathbb{A}}
\newcommand{\PP}{\mathbb{P}}
\newcommand{\E}{\mathbb{E}}
\newcommand{\eps}{\varepsilon}
\newcommand{\dd}{\,\mathrm{d}}
\begin{document}
\section*{UnsolvedMath ID: 701064}
\textbf{Source problem code:} catalogue entry
\subsection{Definitions and mathematical statement}
All groups considered are finite. A Fitting class is a nonempty class of groups, closed under isomorphisms, normal subgroups, and products of normal subgroups belonging to the class: if \(A,B\trianglelefteq X\) and \(A,B\in\mathfrak K\), then \(AB\in\mathfrak K\). A class is soluble if every group belonging to it is soluble.

For a Fitting class \(\mathfrak M\), let \(X_{\mathfrak M}\) denote the largest normal subgroup of \(X\) belonging to \(\mathfrak M\). Under the usual Fitting-product convention, put
\[
\mathfrak M\mathfrak H
=\{X : X/X_{\mathfrak M}\in\mathfrak H\}.
\]
The supplied record does not define this product, so its convention requires confirmation against the cited source. Write \(\operatorname{lFit}(G)\) for the least local Fitting class containing \(G\), when such a least class exists. Thus \(G\in\operatorname{lFit}(G)\), and \(\operatorname{lFit}(G)\subseteq\mathfrak L\) for every local Fitting class \(\mathfrak L\) containing \(G\). The technical term ``local Fitting class'' is retained from the source: its defining local construction is not supplied and requires verification before a fully expanded formalization can be given.

The proposed assertion is the following. For every finite group \(G\) and every pair of Fitting classes \(\mathfrak M,\mathfrak H\), suppose that
\[
\mathfrak F=\mathfrak M\mathfrak H=\operatorname{lFit}(G)
\]
is a soluble, one-generated local Fitting class, and that \(\mathfrak M\ne\mathfrak F\). Must \(\mathfrak H\) be a local Fitting class? Here ``one-generated'' refers to generation by the single group \(G\) as a local Fitting class, not to cyclicity of \(G\). The question remains subject to the unresolved source conventions described above.
\subsection{Short English statement}
If a soluble local Fitting class generated by one finite group is a Fitting product \(\mathfrak M\mathfrak H\), with \(\mathfrak M\) different from the whole class, must \(\mathfrak H\) also be local? The precise local-class convention requires source verification.
\subsection{Sources}
\begin{itemize}
\item Supplied local UnsolvedMath registry record, permanent ID \(701064\), exact source code \texttt{catalogue entry}, titled ``Kourovka Notebook Problem 14.79''; upstream revision \texttt{d371ee917ac9647fbb06cd6fc529deec6f3dda82}. The record was created and updated on 2026-10-16 and supplies the status ``open''; that status is not independently verified here.
\item Attributed in the supplied record to A. N. Skiba, proposed in 1999, in \emph{The Kourovka Notebook (Unsolved Problems in Group Theory)}, 14th issue (1999), Problem 14.79. The supplied background also identifies the 21st issue edition (2026), \texttt{arXiv:1401.0300v48}.
\item The supplied collection description identifies E. I. Khukhro and V. D. Mazurov as editors of the Kourovka Notebook collection covering Issues 1--21 (1965--2026). These bibliographic details are supplied attribution; the cited source has not been independently consulted here.
\end{itemize}
\end{document}
